% 提案システムの目的
This paper proposes a system that estimates the activity levels of office workers using a light detection and ranging (LiDAR) sensor network.
% 点群の特徴抽出
Point-cloud data of individuals at work are captured by LiDAR and converted into numerical features.
% 機械学習による推論
Machine-learning models are employed to estimate the activity levels from the features.
\deleted{Activity level is defined as the degree of concentration during PC-based deskwork.}
\revised{Activity level is defined as a measure of learning efficiency during comprehension of meeting explanations and training content in offices.}
\deleted{Image-based and centroid-based features were extracted from point-cloud data.}
\revised{Image-based and centroid-based features are extracted from point-cloud data.}
% 評価実験
Evaluation experiments were conducted in office environments while participants watched learning videos.
% 活動レベルの評価
Activity indicators were derived from objective evaluation and subjective evaluation, and binary labels were generated using clustering for model evaluation.
% モデルの評価
Several models utilizing different feature types were developed and assessed against the binary labels.
% 結果
\deleted{The results demonstrate that three-dimensional features are effective for activity level estimation.}
\revised{Eight of the ten RF and XGBoost models exceeded a random-inference baseline of 0.5.}
